\section{导读：K2 为什么是“缸中之脑”长出手}

本节先建立整期的技术主线。Kimi K2 是一个基于 MoE 架构的开源编程和 Agentic 大语言模型。访谈里最有力的比喻是“缸中之脑”：只会思考和输出文本的模型像被放在鱼缸里的大脑，和外部世界没有真实连接；而当模型具备编程能力、多轮工具使用能力和外部反馈时，它就开始长出“手”，能操控数字世界。

这期要理解四件事。第一，K2 为什么把基础模型、token efficiency、Agentic 能力和开源放在一起。第二，test-time scaling 为什么同时指向长思考 RL 和 Agent RL。第三，数据墙、FLOPs scaling、Linear Attention、长上下文架构和模型“智商”之间怎样取舍。第四，杨植麟如何把技术判断、商业模式和创始人心态放进“站在无限的开端”这个叙事里。

\begin{figure}[H]
\centering
\includegraphics[width=0.96\textwidth]{brain-in-vat-to-agent.png}
\caption{从缸中之脑到 Agent：编程能力和工具使用让模型开始操控数字世界。自制概念图，依据 00:01:19--00:24:58 对谈内容整理。}
\end{figure}

\begin{knowledgebox}{读图：为什么“手”是关键}
模型会思考和会行动是两件事。Code、工具调用、API、文件系统和反馈环境，让模型的推理可以改变外部状态。Agentic LLM 的关键不只是多说几步，而是多轮使用工具并根据结果修正。
\end{knowledgebox}

\begin{importantbox}{本期核心命题}
K2 的意义不只是“又一个强模型”，而是把基础模型、编程能力、工具使用、开源生态和 test-time scaling 放进同一个系统：让模型从封闭认知系统走向可行动的数字智能体。
\end{importantbox}

\begin{knowledgebox}{本期系统总览}
\begin{tabularx}{\textwidth}{p{0.20\textwidth}p{0.34\textwidth}X}
\toprule
层级 & 核心问题 & 访谈里的答案 \\
\midrule
哲学层 & 为什么继续攀登？ & 问题不可避免，但问题可以解决。 \\
模型层 & K2 怎样成为强底座？ & MoE、token efficiency、Muon、Rephrase。 \\
行动层 & 模型怎样长出手？ & 编程能力、多轮工具使用、Agentic 泛化。 \\
生态层 & 为什么开源？ & 让基础模型能力被外部验证和生态放大。 \\
商业层 & 怎么赚钱、边界在哪？ & API、Agent 产品、基座模型公司和应用公司重新划线。 \\
组织层 & 创始人如何承受复杂性？ & 用 RL 式反馈管理，但警惕被指标 hack。 \\
\bottomrule
\end{tabularx}
\end{knowledgebox}

\begin{knowledgebox}{术语消化：本期关键词}
\begin{tabularx}{\textwidth}{p{0.22\textwidth}p{0.32\textwidth}X}
\toprule
术语 & 含义 & 为什么重要 \\
\midrule
MoE & Mixture of Experts，按 token/任务路由到不同专家 & 提高参数容量，同时控制每次激活计算。 \\
Agentic LLM & 能多轮使用工具、观察反馈并完成任务的语言模型 & 从聊天走向数字世界行动。 \\
Test-time scaling & 推理时投入更多计算、思考或工具调用 & 让模型在测试阶段继续扩展能力。 \\
Token efficiency & 每个训练 token 带来的能力收益 & 数据墙时代，喂同样数据要“脑子长得更多”。 \\
Muon & 访谈中提到的优化器方向 & 可能提升训练效率，但训练稳定性有挑战。 \\
\bottomrule
\end{tabularx}
\end{knowledgebox}

\subsection{本章小结}

EP113 是 K2 发布后的技术判断访谈。它把 Agentic LLM 放在模型公司战略、开源、训练效率、架构取舍和创始人心理之中，适合与 EP119 的 Attention 架构综述、EP115 的 Agent 理论框架对读。

